% !TEX root = task-2-2.tex
% Editable task bank. Append new BankTask blocks at the end of the document.
\documentclass[12pt,a4paper]{article}
\usepackage{physics-textbook}
\usepackage{graphicx,needspace}
\usetikzlibrary{decorations.pathmorphing}
\newif\ifShowSolutions
\ifdefined\TutorEdition
  \ShowSolutionstrue
\else
  \ShowSolutionsfalse % Change to \ShowSolutionstrue to display solutions.
\fi
\newcounter{banktask}
\newsavebox{\BankTaskBox}
\newcommand{\BankTaskContent}[3]{%
  \noindent\textbf{\thebanktask.}\quad #2\par
  \ifShowSolutions
    \Needspace{50mm}% Keep the solution heading with its opening text or diagram.
    #3\par
    \ifdefined\TutorEdition
      \ifcsname TutorSource#1\endcsname
        \nopagebreak\smallskip{\small\csname TutorSource#1\endcsname\par}%
      \fi
    \fi
  \fi
}
\newcommand{\BankTask}[3]{%
  \refstepcounter{banktask}\par\medskip
  \sbox{\BankTaskBox}{\begin{minipage}{\linewidth}%
    \BankTaskContent{#1}{#2}{#3}\end{minipage}}%
  \ifdim\dimexpr\ht\BankTaskBox+\dp\BankTaskBox\relax>\dimexpr\textheight-20mm\relax
    \Needspace{35mm}\BankTaskContent{#1}{#2}{#3}%
  \else
    \Needspace{\dimexpr\ht\BankTaskBox+\dp\BankTaskBox\relax}%
    \noindent\usebox{\BankTaskBox}\par
  \fi
  \medskip
}
\begin{document}
\BookSubsection{2.2}{Модели газов, жидкостей и твёрдых тел}


% Source sdamgia-810: https://phys-ege.sdamgia.ru/problem?id=810
\BankTask{sdamgia-810}{
Какое из утверждений справедливо для кристаллических тел?

$1$) во время плавления температура кристалла изменяется

$2$) в расположении атомов кристалла отсутствует порядок

$3$) атомы кристалла расположены упорядоченно

$4$) атомы свободно перемещаются в пределах кристалла
}{
\smallskip\noindent\textbf{Решение}\par
Кристалл — это агрегатное состояние вещества, которое характеризуется строгим порядком расположения атомов. Атомы в кристаллах не могут свободно перемещаться, колеблясь около узлов кристаллической решетки. Во время плавления температура кристалла не меняется, вся поступающая энергия идет на разрушение кристаллической решетки. Верно утверждение $3$.

Ответ: $3$.
}

% Source sdamgia-827: https://phys-ege.sdamgia.ru/problem?id=827
\BankTask{sdamgia-827}{
Какая-либо упорядоченность в расположении частиц вещества отсутствует. Это утверждение соответствует модели

$1$) только газа

$2$) только жидкости

$3$) только твердого тела

$4$) газа, жидкости и твердого тела
}{
\smallskip\noindent\textbf{Решение}\par
Полное отсутствие какой-либо упорядоченности в расположении частиц вещества соответствует только модели газа.

Ответ: $1$.
}

% Source sdamgia-828: https://phys-ege.sdamgia.ru/problem?id=828
\BankTask{sdamgia-828}{
В процессе перехода вещества из жидкого состояния в кристаллическое

$1$) существенно увеличивается расстояние между его молекулами

$2$) молекулы начинают притягиваться друг к другу

$3$) существенно увеличивается упорядоченность в расположении его молекул

$4$) существенно уменьшается расстояние между его молекулами
}{
\smallskip\noindent\textbf{Решение}\par
В процессе перехода вещества из жидкого состояния в твердое происходит построение кристаллической решетки, а значит, существенно увеличивается упорядоченность в расположении молекул вещества.

Ответ: $3$.
}

% Source sdamgia-836: https://phys-ege.sdamgia.ru/problem?id=836
\BankTask{sdamgia-836}{
Какие частицы находятся в узлах решетки металла?

$1$) нейтральные атомы

$2$) электроны

$3$) отрицательные частицы

$4$) положительные ионы
}{
\smallskip\noindent\textbf{Решение}\par
В узлах решетки металла находятся положительные ионы. Валентные электроны в атомах металлов слабо связаны с ядром. Поэтому в металле они обобществляются, перестают принадлежать определенному атому и образуют так называемый «электронный газ». Атомы металла в узлах решетки, потерявшие свои валентные электроны, превращаются в положительные ионы.

Ответ: $4$.
}

% Source sdamgia-837: https://phys-ege.sdamgia.ru/problem?id=837
\BankTask{sdamgia-837}{
Какое свойство отличает монокристалл от аморфного тела?

$1$) прочность

$2$) электропроводность

$3$) прозрачность

$4$) анизотропность
}{
\smallskip\noindent\textbf{Решение}\par
Аморфное тело не имеет кристаллической структуры и, как правило, — изотропно, то есть не обнаруживают различных свойств в разных направлениях. Монокристаллом же называется отдельный однородный кристалл, имеющий непрерывную кристаллическую решетку и характеризующийся анизотропией свойств. Таким образом, главным свойством, отличающим монокристалл от аморфного тела, является анизотропность.

Ответ: $4$.
}

% Source sdamgia-3205: https://phys-ege.sdamgia.ru/problem?id=3205
\BankTask{sdamgia-3205}{
Какое из утверждений не соответствует представлениям молекулярно-кинетической теории о строении газов?

$1$. Все тела состоят из частиц: атомов, молекул и ионов;

$2$. Частицы находятся в непрерывном хаотическом движении (тепловом);

$3$. Все частицы летают со строго определенными по величине скоростями;

$4$. Частицы взаимодействуют друг с другом путем абсолютно упругих столкновений.
}{
\smallskip\noindent\textbf{Решение}\par
Согласно молекулярно-кинетической теории, все тела, состоят из частиц, которые находятся в непрерывном хаотическом движении и взаимодействуют только посредством упругих столкновений. Естественно, что при столкновениях скорости частиц изменяются. Утверждение $3$ не является верным.

Ответ: $3$.
}

% Source sdamgia-3207: https://phys-ege.sdamgia.ru/problem?id=3207
\BankTask{sdamgia-3207}{
Какое из утверждений справедливо для газа?

$1$. Молекулы газа располагаются в строгом порядке

$2$. Газ сохраняет начальный объем

$3$. Характерное расстояние между молекулами газа по порядку величины совпадает с размерами самих молекул

$4$. Газ всегда занимает весь отведенный ему объем
}{
\smallskip\noindent\textbf{Решение}\par
Из всех приведенных утверждений для газа справедливо только последнее. За счет теплового движения, молекулы газа свободно перемещаются по отведенному им объему. Второе утверждение противоречит четвертому. Первое утверждение относится к кристаллическим телам, третье справедливо для твердых тел и жидкостей.

Ответ: $4$.
}

% Source sdamgia-3208: https://phys-ege.sdamgia.ru/problem?id=3208
\BankTask{sdamgia-3208}{
Какие из утверждений справедливы для жидкостей?

А) Характерное расстояние между молекулами жидкости по порядку величины совпадает с размерами самих молекул 

Б) Жидкость сохраняет форму

В) Жидкость сохраняет объем

Г) Молекулы жидкости образуют периодичную решетку

$1$. А и Г 

$2$. Б и Г

$3$. А и В

$4$. Б и В
}{
\smallskip\noindent\textbf{Решение}\par
Молекулы жидкости достаточно плотно упакованы, так что утверждение А верно. В силу плотной упаковки и достаточно сильного взаимодействия между молекулами, жидкости практически несжимаемы. Как следствие, они сохраняют объем. Молекулы жидкости ведут так называемый «кочевой образ жизни», они могут переходить из одного положения в другое, этим обеспечивается текучесть жидкостей, таким образом, жидкости не сохраняют форму. Утверждение Г относится к кристаллическим телам. Окончательно, верными являются утверждения А и В.

Ответ: $3$.
}

% Source sdamgia-3211: https://phys-ege.sdamgia.ru/problem?id=3211
\BankTask{sdamgia-3211}{
Вещество сохраняет объем, но не сохраняет форму. Это утверждение соответствует модели

$1$) только газа

$2$) только жидкости

$3$) только твердого тела

$4$) газа, жидкости и твердого тела
}{
\smallskip\noindent\textbf{Решение}\par
Твердое тело имеет кристаллическую решетку, поэтому оно сохраняет и форму и объем. Газ занимает весь отведенный ему объем, так что не сохраняет ни форму, ни объем. Молекулы жидкости упакованы практически также плотно, как и молекулы кристаллических тел, и взаимодействуют достаточно сильно, поэтому жидкость сохраняет объем. Однако молекулы жидкости не образуют решетки, они ведут так называемый «кочевой образ жизни», что обеспечивает текучесть жидкости. Как результат жидкости не сохраняют форму.

Ответ: $2$.
}

% Source sdamgia-3299: https://phys-ege.sdamgia.ru/problem?id=3299
\BankTask{sdamgia-3299}{
Расстояние между соседними частицами вещества в среднем во много раз превышает размеры самих частиц. Это утверждение соответствует

$1$) только модели строения газов

$2$) только модели строения аморфных тел

$3$) модели строения газов и жидкостей

$4$) модели строения газов, жидкостей и твердых тел
}{
\smallskip\noindent\textbf{Решение}\par
Утверждение о том, что расстояние между соседними частицами вещества в среднем во много раз превышает размеры самих частиц, соответствует только модели строения газов. В жидком и твердом агрегатных состояниях частицы вещества «упакованы» гораздо плотнее.

Ответ: $1$.
}

% Source sdamgia-3789: https://phys-ege.sdamgia.ru/problem?id=3789
\BankTask{sdamgia-3789}{
Какое из приведенных ниже утверждений справедливо для кристаллических тел?

$1$) в расположении атомов отсутствует порядок

$2$) атомы свободно перемещаются в пределах тела

$3$) при изобарном плавлении температура тела остается постоянной

$4$) при одинаковой температуре диффузия в кристаллах протекает быстрее, чем в газах
}{
\smallskip\noindent\textbf{Решение}\par
Твердые тела обладают кристаллической решеткой, в расположении атомов присутствует дальний порядок. Атомы располагаются в узлах решетки, и в отличие от газов не могут свободно перемещаться по объему тела. Это в частности приводит к тому, что при прочих равных условиях диффузия в твердых телах протекает гораздо медленнее, чем в газах. Таким образом, утверждения $1$, $2$ и $4$ не верны. А вот утверждение $3$ справедливо. Опыт показывает, что при изобарном плавлении, температура твердого тела не изменяется, все поступающее к телу тепло идет на разрушение кристаллической решетки.

Ответ: $3$.
}

% Source sdamgia-4343: https://phys-ege.sdamgia.ru/problem?id=4343
\BankTask{sdamgia-4343}{
Модель, служащая для демонстрации внутреннего строения тел, устроена следующим образом: в дно прямоугольной коробки воткнуты одинаковые упругие вертикальные стерженьки, на каждый из которых насажен магнитик в виде плоской таблетки. После приведения одного из магнитиков в колебательное движение вскоре начинают хаотически колебаться на стерженьках и все остальные магнитики, отталкиваясь друг от друга. Эта модель лучше всего иллюстрирует поведение молекул

$1$) идеального газа

$2$) жидкости

$3$) твердого тела

$4$) идеального газа и жидкости
}{
\smallskip\noindent\textbf{Решение}\par
Данная модель лучше всего демонстрирует поведение молекул твердого тела, потому что именно в данном агрегатном состоянии частицы совершают малые колебания вокруг положений равновесия, не перемещаясь по объему тела.

Ответ: $3$
}

% Source sdamgia-7695: https://phys-ege.sdamgia.ru/problem?id=7695
\BankTask{sdamgia-7695}{
В стакане находится вода в кристаллическом состоянии при температуре $0^\circ$C. Из-за контакта с теплым стаканом и с комнатным воздухом вся вода переходит в жидкое состояние. В процессе перехода воды из кристаллического состояния в жидкое

$1$) ее объем уменьшается;

$2$) ее плотность уменьшается;

$3$) ее внутренняя энергия уменьшается;

$4$) ее температура уменьшается.
}{
\smallskip\noindent\textbf{Решение}\par
При переходе воды из кристаллического состояния в жидкое, ее плотность увеличивается, а значит при неизменной массе вещества объем должен уменьшиться.

Ответ: $1$
}

% Source sdamgia-3464: https://phys-ege.sdamgia.ru/problem?id=3464
\BankTask{sdamgia-3464}{
В молекулярной физике используется понятие «идеальный газ». Это понятие применимо тогда, когда можно пренебречь:

А. Средней кинетической энергией поступательного движения атомов или молекул.

Б. Энергией взаимодействия атомов или молекул, то есть действием сил притяжения или отталкивания.

В. Массой атомов или молекул.

Какое(-ие) из утверждений правильно(-ы)?

$1$) только А

$2$) только Б

$3$) только В

$4$) А и В
}{
\smallskip\noindent\textbf{Решение}\par
Идеальный газ — это математическая модель газа, в которой предполагается, что потенциальной энергией взаимодействия молекул (или атомов) можно пренебречь по сравнению с их кинетической энергией. В этой модели предполагается, что между молекулами не действуют силы притяжения или отталкивания, соударения частиц между собой и со стенками сосуда абсолютно упруги.

Таким образом, утверждение Б верно, утверждение А неверно. Поскольку кинетической энергией в модели идеального газа не пренебрегают, то не могут пренебрегать и массой частиц газа, так как кинетическая энергия этой массе пропорциональна. Следовательно, для применимости модели идеального газа необходимо выполнение только требования Б.

Ответ: $2$.
}

% Source sdamgia-3543: https://phys-ege.sdamgia.ru/problem?id=3543
\BankTask{sdamgia-3543}{
Какие положения включает модель идеального газа?

А. Объемом, занимаемым молекулами газа, можно пренебречь по сравнению с объемом сосуда.

Б. Взаимодействием молекул можно пренебречь.

В. Молекулы не сталкиваются друг с другом при движении.

$1$. А, Б и В

$2$. А и Б

$3$. А и В

$4$. Б и В
}{
\smallskip\noindent\textbf{Решение}\par
Идеальный газ — математическая модель газа, в которой предполагается, что потенциальной энергией взаимодействия молекул можно пренебречь по сравнению с их кинетической энергией. В этой модели считается, что между молекулами не действуют силы притяжения или отталкивания, то есть взаимодействием молекул пренебрегают. Модель идеального газа хорошо работает только для не очень плотных газов, поскольку в этой модели пренебрегают объемом, занимаемым молекулами газа по сравнению с объемом сосуда. А вот утверждение В к этой модели не относится. Поскольку именно за счет столкновений молекул, в результате которых они обмениваются кинетическими энергиями, газ приходит к состоянию термодинамического равновесия.

Ответ: $2$.
}

% Source sdamgia-4118: https://phys-ege.sdamgia.ru/problem?id=4118
\BankTask{sdamgia-4118}{
В учебнике по физике в одном из абзацев написано: «Молекулы считаются материальными точками, которые хаотически движутся и абсолютно упруго соударяются друг с другом и со стенками сосуда. В промежутках между столкновениями молекулы друг с другом и со стенками сосуда не взаимодействуют». Какая физическая модель описывается в этом абзаце учебника?

$1$) монокристаллическое твердое тело

$2$) поликристаллическое твердое тело

$3$) идеальная жидкость

$4$) идеальный газ
}{
\smallskip\noindent\textbf{Решение}\par
Данное описание относится к модели идеального газа.

Ответ: $4$.
}

\end{document}
